\documentclass[11pt]{article}
\usepackage[T1]{fontenc}
\usepackage{textcomp}
\usepackage[margin=0.75in]{geometry}
\usepackage{amsmath,amssymb,amsthm}
\usepackage{array,booktabs}
\usepackage{hyperref}
\setlength{\parindent}{0pt}
\newtheorem{theorem}{Theorem}
\theoremstyle{definition}
\newtheorem{definition}{Definition}
\title{Flow Proof Artifact: prop-21-ex-aequali-three-ratios}
\author{Generated by \texttt{flow doc proof}}
\date{Flow Proof Book}
\begin{document}
\maketitle
If a first is to a second as a third is to a fourth, and a fifth is to a sixth as the third is to a fourth, the sum of the first and fifth is to the second as the sum of the third and seventh is to the fourth.\\[0.5em]
\textbf{Source.} euclid — Elements, Book V, Proposition 21\\[0.5em]
\bigskip
\noindent{\large \textbf{Derived fact 1.} \textit{If a first is to a second as a third is to a fourth, and a fifth is to a sixth as the third is to a fourth, the sum of the first and fifth is to the second as the sum of the third and seventh is to the fourth} \hfill the Euclidean plane \cdot Euclid Book V \cdot Proposition 21: ex aequali from three linked ratios}\par
\smallskip\noindent\textit{Coordinate: the Euclidean plane \cdot Euclid Book V \cdot Proposition 21: ex aequali from three linked ratios}\par
\smallskip\noindent\textit{Source: euclid — Elements, Book V, Proposition 21}\par
\smallskip\noindent\textit{Built on: proposition 11: proportional magnitudes satisfy componendo, for Euclid Book V on the Euclidean plane, proposition 19: ex aequali proportion in the second form, for Euclid Book V on the Euclidean plane}\par
\medskip
\noindent\textbf{Goal.} If a first is to a second as a third is to a fourth, and a fifth is to a sixth as the third is to a fourth, the sum of the first and fifth is to the second as the sum of the third and seventh is to the fourth\par
\noindent$\displaystyle \text{ex aequali three ratios holds}$\par
\medskip
\noindent
\begin{tabular}{@{} >{\raggedright\arraybackslash}p{0.06\textwidth} >{\raggedright\arraybackslash}p{0.47\textwidth} >{\raggedleft\arraybackslash}p{0.41\textwidth} @{}}
\textbf{\#} & \textbf{Proof} & \textbf{Mathematics} \\
\midrule
\textbf{1.} & We prove that proposition 21: ex aequali from three linked ratios for Euclid Book V the Euclidean plane. & \\[0.45em]
\textbf{2.} & Let a = first. & \\[0.45em]
\textbf{3.} & Let b = second. & \\[0.45em]
\textbf{4.} & Let e = fifth. & \\[0.45em]
\textbf{5.} & Let g = seventh. & \\[0.45em]
\textbf{6.} & From step 2, step 3, step 4, and step 5, we can deduce that ratio(a plus e, b) equals ratio(c plus g, d). & $\displaystyle ratio(a + e, b) = ratio(c + g, d)$ \\[0.45em]
\textbf{7.} & We invoke the derived fact governing Euclid Book V the Euclidean plane: proposition 11: proportional magnitudes satisfy componendo, for Euclid Book V on the Euclidean plane. & \\[0.45em]
\textbf{8.} & We invoke the derived fact governing Euclid Book V the Euclidean plane: proposition 19: ex aequali proportion in the second form, for Euclid Book V on the Euclidean plane. & \\[0.45em]
\textbf{9.} & From step 7 and step 8, this implies ex aequali three ratios holds. Hence proven. & $\displaystyle \text{ex aequali three ratios holds}$ \\[0.45em]
\bottomrule
\end{tabular}
\label{thm:Geometry-Euclid Book V-proposition_21_ex_aequali_from_three_linked_ratios}
\par\medskip\hrule
\end{document}
